% ============================================================
% ch21b.tex —— 第 21 章「S 类理论」后半（21.5 起至章末）
% 底本：LNM 2386，PDF p383–p404（印刷 pp.343–364）
% 续块：不写 \chapter
% ============================================================

% [原文 p.343]
\section{$A_1$ 型 S 类理论}

\begin{definition}[21.24] 若 Hitchin 场在某穿孔处的极点阶数（至多）为一，则称该穿孔为正则的；若极点阶数更高，则称该穿孔为非正则的。
\end{definition}

\subsection{广义 SU(2) 箭图}

现在转向更一般的、涉及 SU(2) 规范群与味群的线性共形箭图（图 21.11）。$\SL_2(\Z)$ 对偶作用在相关 $4d$ $\mathcal{N}=2$ 理论的弱耦合描述上，可通过聚焦于除一个规范群之外其余规范群都极弱耦合的情形来理解。剩下的那个规范群表现得像带四个基本味的 SU(2) SYM 理论。把这最后一个群的规范耦合增大到它变为强耦合，它就可以更好地用一个对偶的弱耦合规范群来描述。

\Fig[0.95\linewidth]{figures/fig_21_11.png}{21.11}{更复杂的 SU(2) 箭图。}

图 21.11 左边的箭图表示一个规范群为 $\SU(2)_1\times\SU(2)_2$ 的 $4d$ $\mathcal{N}=2$ 理论，其中每个规范因子有两个基本味。此外还有一个处于双基本表示 $\mathbf{2}_1\otimes\mathbf{2}_2$ 的超多重态。这个双基本超多重态引入一个额外的 SU(2) 味对称性\footnote{该对称性是 $\SU(2)\simeq\mathrm{USp}(2)$ 而非 $\U(1)$，其原因在于 $\mathbf{2}_1\otimes\mathbf{2}_2$ 是实表示。}，图中记为 $\SU(2)_e$。因此每个 SU(2) 规范因子实际上有四个基本味，而规范耦合 $\tau_1$ 与 $\tau_2$ 是该理论的精确边缘参数。该理论的 S 对偶性质最早在 [23] 中研究。

我们先考虑两个规范群都弱耦合的对偶框架，并增大 $\SU(2)_1$ 的规范耦合 $g_1$、同时保持 $g_2$ 很小。于是规范群 $\SU(2)_1$ 更好地被描述为另一个弱耦合的 SU(2) 规范因子。如前所述，这种描述的改变一般会涉及
% [原文 p.344]
与 $\SU(2)_1$ 相关联的「味」对称因子的所有可能置换，即 $\SU(2)_{a,b,e}$ 与 $\SU(2)_2$。当交换两个规范因子 $\SU(2)_1$ 与 $\SU(2)_2$ 的角色时，同样的推理也适用。于是，通过对两个规范因子之一反复作用对偶群 $\SL_2(\Z)$，味对称因子 $\SU(2)_{a,b,c,d,e}$ 会以所有可能的方式被交换。重要的是，人们看到与这两个规范因子相关联的两份 $\SL_2(\Z)$ 并不交换。

该理论的 Seiberg–Witten 曲线可用 21.3 节讨论的方法计算，并且仍可写成如下形式：
\begin{equation}
\lambda^2-\phi_2(z)=0, \tag{21.15}
\end{equation}
其中 $\lambda=x\,\dd z/z$ 是 Seiberg–Witten 微分，$\phi_2(z)$ 是球面 $\mathbb{CP}^1_z$ 上带五个简单穿孔的二次微分。相应地：

\begin{proposition}[21.25] 该理论的精确边缘形变空间可等同于带五个穿孔的球面的泰希米勒空间 $T(S_{0,5})$。映射类群 $\mathrm{Mod}(S_{0,5})$ 作为对偶群作用。模空间 $T(S_{0,5})/\mathrm{Mod}(S_{0,5})$ 有一个尖点，它对应于该理论形如图 21.11 左图那种箭图的弱耦合描述。限制到 $\mathrm{Mod}(S_{0,5})$ 中固定穿孔的子群 $\widetilde{\mathrm{Mod}}(S_{0,5})$，商 $T(S_{0,5})/\widetilde{\mathrm{Mod}}(S_{0,5})$ 有 15 个尖点，对应于味因子 $\SU(2)_{a,b,c,d,e}$ 的所有可能配置。几何上，这些尖点中的每一个都对应于把 $S_{0,5}$ 退化为由细长管道连接的三个球面：其中两个球面各带两个穿孔、各连一根管道，第三个球面带一个穿孔并连两根管道。该理论的 UV 曲线是带五个穿孔的 $\mathbb{CP}^1$。
\end{proposition}

现在转向图 21.11 右边的箭图。在第一个与第三个规范因子处对箭图做对偶，会得到一个外观相似的箭图，只是味因子被打乱了。在第二个规范节点处对箭图做对偶，则给出 $\SU(2)_2$ 的四个「味」因子 $\SU(2)_{1,c,d,3}$ 分成两对的所有可能方式。

其中一种安排给出图 21.12 所示的箭图，它是一个广义 SU(2) 箭图的例子。它所描述的弱耦合理论可以看作由图 21.11 左图得到：把味对称因子 $\SU(2)_e$ 规范化，并为这个新的规范群添加两个基本味，以保证标度不变性。

同样，该理论的 SW 曲线可写成与式 (21.15) 相同的形式，其中 $\phi_2$ 是 $\mathbb{CP}^1_z$ 上带六个简单穿孔的二次微分。该理论作为 SU(2) 箭图的弱耦合描述，对应于图 21.13 所画的 $\mathbb{CP}^1_z$ 带六个简单穿孔的退化。

精确边缘形变的空间可等同于带六个穿孔的球面的泰希米勒空间，映射类群在其上作为对偶群作用。

图 21.11 的箭图是涉及 SU(2) 规范群与味群的线性共形箭图的具体情形。通过对规范节点作所有可能的对偶变换，
% [原文 p.345]
\Fig[0.4\linewidth]{figures/fig_21_12.png}{21.12}{一个广义 SU(2) 箭图。}
\Fig[0.8\linewidth]{figures/fig_21_13.png}{21.13}{带六个穿孔的球面的可能退化。}
得到的箭图不一定是线性的，但总是对应于带外边的三叉树。在这种表示中，每条内部边表示一个 SU(2) 规范因子，每条外部边表示一个 SU(2) 味因子，而每个三叉顶点对应于四个超多重态的一种配置，具体是下列之一：
\begin{itemize}
\item 某个 SU(2) 规范因子的一对基本超多重态。味对称性为 $\SU(2)\times\SU(2)$。
\item 处于两个 SU(2) 规范因子的双基本表示中的一个超多重态，相应的味对称性为 SU(2)。
\item 处于三个 SU(2) 规范因子的三基本表示中的一个超多重态，相应的味对称性平凡。
\end{itemize}
这些图是连通的三叉树，这意味着若 $E_i=n$ 是内部边的数目，则必有 $V=n+1$ 个三叉顶点，以及
\[
E_e=3V-2E_i=3(n+1)-2n=n+3
\]
条外部边。箭图描述使味对称群的一个 $\SU(2)^{n+3}$ 子群显式可见。借助 S 对偶，可以达到具有 $n$ 条内部边的所有可能三叉树拓扑。对固定的 $n$，每一种这样的拓扑都对应于某个唯一底层理论的一种特定弱耦合描述。等价地说，
% [原文 p.346]
这些描述对应于带 $n+3$ 个穿孔的黎曼球面的不同退化极限。这类理论的 Seiberg–Witten 曲线可写成式 (21.15) 的形式，其中 $\phi_2(z)$ 是 $\mathbb{CP}^1_z$ 上带 $n+3$ 个简单穿孔的二次微分。该理论的精确边缘形变空间等同于带 $n+3$ 个穿孔的球面的泰希米勒空间 $T(S_{0,n+3})$，因此这个唯一的底层理论自然地与 $S_{0,n+3}$ 相联系。

\begin{remark}[21.26] 弱耦合描述与 $S_{0,n+3}$ 的裤子分解一一对应（参见 4.4 节）。
\end{remark}

还可以进一步推广所考虑的箭图，办法是放宽它们是树这一限制。特别地，Witten 对 SU(2) 椭圆模型（elliptic models）的构造给出项链箭图。对这种箭图应用 S 对偶，可得带外边的所有可能的三叉图（带一个回路）。图 21.14 左图即为一例。在这种情况下，Seiberg–Witten 曲线可写成
\[
\lambda^2-\phi_2(z)=0,
\]
其中 $\lambda=x\,\dd z$ 是椭圆曲线 $E_z$ 的全纯余切丛上的 Liouville 1-形式，限制到 Seiberg–Witten 曲线上时与 Seiberg–Witten 微分等同，而 $\phi_2(z)$ 是 $E_z$ 上带简单穿孔的二次微分。例如，图 21.14 左图所示的箭图对应于 $E_\tau$ 上带 6 个简单穿孔的二次微分 $\phi_2$。这里该理论的 UV 曲线是带 6 个穿孔的椭圆曲线 $E_\tau$。

具有固定内部边（即 SU(2) 规范因子）数目 $n$ 的这族单回路箭图，编码了与带 $n$ 个穿孔的光滑环面 $S_{1,n}$ 相关联的某个底层理论的所有可能弱耦合描述。等价地，这类箭图与 $S_{1,n}$ 到裤子的不等价分解一一对应。

更一般地，可以考虑带 $g$ 个回路的广义 SU(2) 箭图。对于这些理论，没有直接的办法借助 IIA 型超弦理论与 M 理论、按 21.3 节的方法构造 Seiberg–Witten 曲线；但线性箭图与项链箭图的例子提示如下。若所考虑的箭图有 $g$ 个回路与 $n$ 条外部边，则 Seiberg–Witten 曲线写成：
\[
\lambda^2-\phi_2(z)=0,
\]
其中 $\lambda=x\,\dd z$ 是亏格 $g$、带 $n$ 个穿孔的黎曼曲面 $C_{g,n}$（UV 曲线）的全纯余切丛上的 Liouville 1-形式，限制到 Seiberg–Witten 曲线上时与 Seiberg–Witten 微分等同，而 $\phi_2(z)$ 是 $C_{g,n}$ 上带 $n$ 个简单穿孔的二次微分。图 21.14 右图展示了两个例子。

如前所述，可以证明：对具有 $g$ 个回路与 $n$ 条外部边的广义 SU(2) 箭图，在其规范节点处反复施加 S 对偶，可以得到具有 $g$ 个回路与 $n$ 条外部边的广义 SU(2) 箭图的所有可能配置。这一过程提示存在一个与亏格 $g$、带 $n$ 个穿孔的光滑曲面 $S_{g,n}$ 相关联的底层 $4d$ $\mathcal{N}=2$ SCFT。该理论
% [原文 p.347]
\Fig[0.85\linewidth]{figures/fig_21_14.png}{21.14}{对应于 $\mathcal{S}[A_1,S_{g,n}]$ 的弱耦合描述的广义 SU(2) 箭图。右图两个箭图被画成带外边的三叉树。}
的精确边缘形变空间等同于泰希米勒空间 $T(S_{g,n})$，映射类群 $\mathrm{Mod}(S_{g,n})$ 在其上作为对偶群作用。该理论的弱耦合描述、亦即相应的广义 SU(2) 箭图，对应于 $S_{g,n}$ 的裤子分解。

\begin{remark}[21.27] 对 $S_{g,n}$ 的每个裤子分解，都存在 $T(S_{g,n})$ 上相应的一组 Fenchel–Nielsen 坐标（见第 7 章）。与该分解相联系的弱耦合极限，对应于 $T(S_{g,n})$ 中那些在裤子分解里出现的 $S_{g,n}$ 上简单闭曲线的长度非常小的区域。在此极限下，Fenchel–Nielsen 坐标可与 UV 理论的参数相联系：每条简单闭曲线的长度对应于相应 SU(2) 规范因子的 UV 规范耦合，而 twist 参数则与其 theta 角 $\Theta$ 等同 [150]。
\end{remark}

\begin{remark}[21.28] 给定 $S_{g,n}$ 上的一条简单闭曲线 $\gamma$，可以考虑其长度 $l(\gamma)$ 趋于零的极限。在此极限下，对任何包含 $\gamma$ 的 $S_{g,n}$ 裤子分解，对应于 $\gamma$ 的那个 SU(2) 规范群会完全解耦。于是：
\begin{itemize}
\item 若 $\gamma$ 是非分离的\footnote{光滑曲面 $S$ 上的一条简单闭曲线 $\gamma$ 称为分离的，若 $S\setminus\gamma$ 有两个连通分支；否则称为非分离的。}，则留下一个对应于 $S_{g-1,n+2}$ 的广义箭图。
\item 若 $\gamma$ 是分离的，则得到两个分别对应于 $S_{g_1,n_1}$ 与 $S_{g_2,n_2}$ 的广义箭图，其中 $g_1+g_2=g$ 且 $n_1+n_2=n+2$。
\end{itemize}
\end{remark}

% [原文 p.348]
这对应于泰希米勒空间 $T(S_{g,n})$（以及相应的模空间）中的无穷远，分别对应于：若 $\gamma$ 非分离则为 $T(S_{g-1,n+2})$，若 $\gamma$ 分离则为乘积 $T(S_{g_1,n_1})\times T(S_{g_2,n_2})$。

通过考虑与 $S_{g,n}$、$S_{g_1,n_1}$，或与 $S_{g_1,n_1}$ 和 $S_{g_2,n_2}$ 相关联的 $4d$ $\mathcal{N}=2$ 理论，更具体地考虑其模空间的分支（库仑分支或希格斯分支），这暗示存在有趣的、取值于全纯辛流形范畴的 $1+1$ 维拓扑量子场论 [272]。

\subsection{六维定义}

在 Witten 对 Seiberg–Witten 曲线的构造 [366] 中，把 IIA 型膜构型提升到 M 理论，通常得到一张在 $\C\times C_{\mathrm{UV}}$ 内缠绕一条复曲线 $\Sigma$ 的 M5 膜，其中 $C_{\mathrm{UV}}$ 表示 UV 曲线——对线性箭图是带穿孔的球面，对项链箭图是带穿孔的椭圆曲线。黎曼曲面 $\Sigma$ 与 D4 膜上世界体积理论的 Seiberg–Witten 曲线等同，并且是 $C_{\mathrm{UV}}$ 上的 2 次分歧覆盖（对 SU(2) 箭图而言）。

缠绕 $\Sigma$ 的这张 M5 膜是如下构型的可归一化形变：两张 M5 膜缠绕 $C_{\mathrm{UV}}$，并与沿 $\C$ 延伸、位于 $C_{\mathrm{UV}}$ 中固定点 $t_a$ 处的横向 M5 膜相交。一对 M5 膜上的世界体积理论是 $A_1$ 型的六维 $\mathcal{N}=(2,0)$ 超共形理论。在足够低的能量下，横向 M5 膜表现为 $6d$ $\mathcal{N}=(2,0)$ 理论中位于 $t_a\in C_{\mathrm{UV}}$ 的余维 2 缺陷。导致单张 M5 膜缠绕 $\Sigma$ 的形变由 $6d$ $\mathcal{N}=(2,0)$ $A_1$ 理论中标量场的真空期望值决定，而这些标量场描述 M5 膜沿 $\C$ 的运动。因此，$4d$ $\mathcal{N}=2$ 的线性或项链共形箭图理论可以定义为 $6d$ $\mathcal{N}=(2,0)$ $A_1$ 理论在 $C_{\mathrm{UV}}$ 上紧化、并带有位于 $t_a\in C_{\mathrm{UV}}$ 的余维 2 简单缺陷时的低能极限。

这一框架可以推广到 $\C^*$ 或 $E_\tau$ 之外的 UV 曲线。给定亏格为 $g$、带 $n$ 个穿孔的黎曼曲面 $C_{g,n}$，可以把 $6d$ $\mathcal{N}=(2,0)$ $A_1$ 理论在 $C_{g,n}$ 上紧化、在穿孔处放置余维 2 简单缺陷，然后取低能极限。这种紧化定义一个 $4d$ $\mathcal{N}=2$ 超共形场论 [146]。为在四维中保持 $\mathcal{N}=2$ 超对称，需要把 $6d$ 理论沿 $C_{g,n}$ 作部分拓扑扭曲（partial twist）。扭曲的效应是：所得的 $4d$ $\mathcal{N}=2$ 理论只依赖于 $C_{g,n}$ 上度量的共形类，即依赖于其复结构。黎曼曲面 $C_{g,n}$ 即为该理论的 UV 曲线。$A_1$ 型的 $6d$ $\mathcal{N}=(2,0)$ 理论有一个 $R$-荷为 2 的保护算子，在扭曲下变为二次微分 $\phi_2(z)$。

所得理论的精确边缘形变空间等同于泰希米勒空间 $T(S_{g,n})$，其中 $S_{g,n}$ 是 $C_{g,n}$ 底下的光滑曲面。对偶群是映射类群 $\mathrm{Mod}(S_{g,n})$。于是，对每个光滑曲面 $S_{g,n}$ 都得到一族理论（模去精确边缘
% [原文 p.349]
形变）。$S_{g,n}$ 上复结构的每种选取——等价地，$T(S_{g,n})$ 中的每个点——对应于把该理论的精确边缘参数取为特定值。

\begin{definition}[21.29] 以这种方式得到的 $4d$ $\mathcal{N}=2$ 理论称为 $A_1$ 型 S 类理论，记作 $\mathcal{S}[A_1,C_{g,n}]$，其中 $C_{g,n}$ 表示 UV 曲线。这类理论的库仑分支对应于 $C_{g,n}$ 上二次微分 $\phi_2(z)$ 构成的空间；在无质量时 $\phi_2$ 在穿孔处有 1 阶极点，有质量时则有 2 阶极点。Seiberg–Witten（SW）曲线 $\Sigma$ 是 $C_{g,n}$ 在 $T^*_{1,0}C_{g,n}$ 中的一个 2 次分歧覆盖，由方程
\[
\lambda^2=\phi_2(z)
\]
确定，其中 $\lambda$ 表示 $T^*_{1,0}C_{g,n}$ 上的全纯 Liouville 1-形式。该理论的精确边缘耦合空间是 $T(S_{g,n})$，其对偶群是 $\mathrm{Mod}(S_{g,n})$。
\end{definition}

\begin{remark}[21.30] 正如上文在带 $F\le3$ 个味的 SU(2) SYM 理论中所指出的，也可以在存在非正则穿孔时紧化 $A_1$ 型的 $6d$ $\mathcal{N}=(2,0)$ 理论。这一进路可以把一大类 $4d$ $\mathcal{N}=2$ 渐近自由理论描述为 S 类理论 [155]。一般地，这些理论记作 $\mathcal{S}[A_1,C_{g,n},\mathcal{D}]$，其中 $\mathcal{D}$ 编码每个穿孔的性质，称为缺陷数据（defect data）。
\end{remark}

\section{$A_2$ 型 S 类理论}

四维 $\mathcal{N}=2$ SU($N$) SYM 带 $F=2N$ 个基本超多重态，不仅对 $N=2$ 是超共形的，对 $N\ge3$ 也是。这些理论的 Seiberg–Witten 曲线最早在 [28] 中算出。值得注意的是，这些曲线不是 $\SL_2(\Z)$ 不变的\footnote{在 $4d$ $\mathcal{N}=2$ SU(2) SYM 带 $F=4$ 个无质量基本超多重态的情形，低能规范耦合 $\tau$ 与 UV 耦合 $\tilde\tau$ 由 $\tau=2\tilde\tau=\theta/\pi+2\pi\ii/g^2$ 相联系。对偶群通过分式线性变换作用在 $\tau$ 上，因此 $T$ 对应于 $\theta\to\theta+\pi$；与 $T^2$ 不同，这并非弱耦合描述的先验对称性。注意 $T$ 以非平凡的方式置换 $\SO(8)$ 的三个 8 维表示。}，而是在同余子群
\[
\Gamma_0(2)=\left\{\begin{pmatrix}a&b\\c&d\end{pmatrix}\in\PSL_2(\Z)\ \middle|\ c\equiv0\ [2]\right\}<\PSL_2(\Z)
\]
下不变。$\HH$ 中 $\PSL_2(\Z)$ 与 $\Gamma_0(2)$ 的基本胞腔连同其 2 阶（紫色菱形）与 3 阶（蓝色三角形）轨形点如图 21.15 所示。

图 21.15 左侧所画的 $\PSL_2(\Z)$ 基本域表明：在相差对偶的意义下，$4d$ $\mathcal{N}=2$ SCFT 总可以描述为带四个基本味、规范耦合满足 $|\tau|\ge1$ 的 SU(2) 规范理论。相应地，
% [原文 p.350]
\Fig[0.95\linewidth]{figures/fig_21_15.png}{21.15}{$\HH$ 中 $\PSL_2(\Z)$ 与 $\Gamma_0(2)$ 的基本胞腔。}
$\HH/\PSL_2(\Z)$ 只有一个尖点。相反，在图 21.15 右侧所示的 $\Gamma_0(2)$ 基本域内，可以取极限 $\tau\to1$。由于 $\HH/\Gamma_0(2)$ 有两个尖点，这一极限先验地不等价于 $\tau\to\infty$。这就提出一个问题：$4d$ $\mathcal{N}=2$ SU($N$) 带 $F=2N$ 的 SYM 如何能以 $\tau\to1$ 来刻画？

\subsection{Argyres–Seiberg 对偶}

带 $F=6$ 个基本超多重态的四维 $\mathcal{N}=2$ SU(3) SYM 的 S 对偶性质由 Argyres 与 Seiberg 在 [29] 中研究。该理论的 Seiberg–Witten 曲线亏格为二；在尖点 $\operatorname{Im}\tau\to\infty$ 处，三个不相交的圈收缩。在此极限下规范耦合 $g$ 趋于 0，完整的 SU(3) 规范对称性恢复。收缩的三个圈对应于 $g=0$ 时三个 $W$ 玻色子（连同其反粒子）变为无质量。相反，当 $\tau\to1$ 时，Seiberg–Witten 曲线上只有一个圈收缩，这表明单个 BPS 粒子变为无质量。因此 $\tau=1$ 处的物理不可能是带 $F=6$ 个基本超多重态的弱耦合 $4d$ $\mathcal{N}=2$ SU(3) SYM 的物理。

单个圈的收缩反而提示一种涉及某个 SU(2) 规范因子的描述，该规范因子在 $\tau\to1$ 时变为弱耦合。然而这种描述必然是不完整的，因为 $4d$ $\mathcal{N}=2$ SU(2) 规范理论的秩为一，而带 $F=6$ 个基本超多重态的 $4d$ $\mathcal{N}=2$ SU(3) SYM 的秩为二。

\begin{proposition}[21.31] 在尖点 $\tau\to1$ 附近，带 $F=6$ 个无质量基本超多重态的 SU(3) SYM 更好地描述为一个弱耦合的 SU(2) 规范理论，它与一个无质量超多重态以及一个孤立的秩 1 $4d$ $\mathcal{N}=2$ SCFT——具体地说是 $E_6$ Minahan–Nemeschansky SCFT [266]——相耦合。
\end{proposition}

SU(2) 规范群与 $E_6$ SCFT 之间的耦合来自对 Minahan–Nemeschansky 理论中整体对称群 $E_6$ 的极大子群 $\SU(2)\times\SU(6)$ 的 SU(2) 因子作规范化。规范化这个 SU(2) 因子把整体对称性 $E_6$ 破缺为 SU(6)，而 SU(2) 基本表示中的赝实超多重态
% [原文 p.351]
贡献一个 $\SO(2)\simeq\U(1)$ 因子。因此总味对称性为 $\SU(6)\times\U(1)\simeq\U(6)$，这与带 $F=6$ 个无质量基本超多重态的 $4d$ $\mathcal{N}=2$ SU(3) SYM 的味对称性相符\footnote{SU(3) 的基本表示是真正的复表示。}。SU(3) 带 $F=6$ 的 SYM 与带一个基本味并与秩 1 $E_6$ SCFT 耦合的 SU(2) SYM 之间的 S 对偶称为 Argyres–Seiberg 对偶。

带 $F=6$ 个无质量基本超多重态的四维 $\mathcal{N}=2$ SU(3) SYM 可以用图 21.16 左边的箭图表示，其中标号 1（相应地 3）的味节点对应于 U(1)（相应地 SU(3)）因子。这一表示来自把六个基本超多重态的集合平均分成两组，对应于味对称群 U(6) 的极大子群 $\U(3)\times\U(3)$。每个 U(3) 因子再分解为 $\U(3)\simeq\SU(3)\times\U(1)$。

Argyres–Seiberg S 对偶理论由图 21.16 右边的箭图表示，其中聚焦于 Minahan–Nemeschansky 理论中 $E_6$ 整体对称性的极大子群 $\SU(3)\times\SU(3)\times\SU(3)$，它以箭图最右端的三叉结点出现。更多细节请参见 [146, 第 3 节]。注意，规范化了 SU(2) 因子的 $E_6$ 理论在这类广义 SU(3) 箭图中所起的作用，类似于广义 SU(2) 箭图中四个超多重态的块所起的作用。

\Fig[0.95\linewidth]{figures/fig_21_16.png}{21.16}{带 $F=6$ 个基本味的 SU(3) SYM 的箭图以及 Argyres–Seiberg 对偶框架。}

图 21.16 的箭图提示：带 $F=6$ 个无质量基本超多重态的四维 $\mathcal{N}=2$ SU(3) SYM 可以与带四个不同类型的穿孔的黎曼球面相联系（同样称为 UV 曲线），如图 21.17 所示，该图画出曲面的两个退化极限。每个退化极限对应于图 21.16 中的一个箭图。在 Argyres–Seiberg 对偶中，对应于 SU(2) 的单个基本超多重态的 U(1) 味因子作为原描述中味因子 $\U(1)_a$ 与 $\U(1)_c$ 的一个组合出现，而另一个组合则是 SU(2) 在 SU(3) 中的中心化子。

这一几何图像的依据同样来自对带 6 个基本超多重态的 SU(3) SYM 的 Seiberg–Witten 曲线的计算，所用方法见 21.3 节。可以把 Seiberg–Witten 曲线（不含质量形变）写成：
\[
(\Sigma):\quad
\lambda^3-\phi_2(z)\lambda-\phi_3(z)
=x^3\frac{\dd z^3}{z^3}
-\frac{u^{(2)}\dd z^2}{(z-1)(z-q)z}\left(x\frac{\dd z}{z}\right)
+\frac{u^{(3)}\dd z^3}{(z-1)(z-q)z^2}=0,
\]

% [原文 p.352]
\Fig[0.9\linewidth]{figures/fig_21_17.png}{21.17}{带 $2+2$ 个穿孔的球面的退化，其中无圈的圆点对应 U(1) 味因子（简单穿孔），带圈的圆点对应 SU(3) 味因子（满穿孔）。}

其中 $\lambda=x\,\dd z/z$ 是 $\mathbb{CP}^1_z$ 的全纯余切丛上的 Liouville 1-形式。注意 $\phi_2(z)$ 是 $\mathbb{CP}^1_z$ 上在 $0,q,1,\infty$ 处有 1 阶极点的二次微分，而 $\phi_3(z)$ 是 $\mathbb{CP}^1_z$ 上在 $z=q,1$ 处有 1 阶极点、在 $z=0,\infty$ 处有 2 阶极点的三次微分，其中 $u^{(2)}$ 与 $u^{(3)}$ 是复参数。

\begin{definition}[21.32] 在无质量形变时，$\mathbb{CP}^1_z$ 上使 $\phi_2$ 有 1 阶极点、$\phi_3$ 有 2 阶极点的点称为满穿孔（full punctures）。使 $\phi_2$ 与 $\phi_3$ 都有 1 阶极点的点称为简单穿孔（simple punctures）。
\end{definition}

在图 21.17 中，简单穿孔 $a$ 与 $c$ 位于 $z=q,1$，而满穿孔 $b$ 与 $d$ 位于 $z=0,\infty$。极限 $\tau\to1$ 对应于让穿孔 $a$ 与 $c$ 相撞，即 $q\to1$。在此极限下，Seiberg–Witten 曲线简化为：
\[
\lambda^3=\frac{u^{(2)}\dd z^2}{(z-1)^2z}\left(x\frac{\dd z}{z}\right)+\frac{u^{(3)}\dd z^3}{(z-1)^2z^2}.
\]
在此极限下，参数 $u^{(2)}$ 可等同于 Argyres–Seiberg 对偶框架中出现的弱耦合 SU(2) 规范群的库仑分支参数。令 $u^{(2)}=0$，曲线化为：
\[
\lambda^3=\frac{u^{(3)}}{(z-1)^2z^2},
\]
它描述 $E_6$ Minahan–Nemeschansky SCFT 的 Seiberg–Witten 曲线。

\begin{remark}[21.33] 当超多重态带有一般质量时，$\phi_2$（相应地 $\phi_3$）在 $z=0,q,1,\infty$ 处有 2 阶（相应地 3 阶）极点。$\lambda|_\Sigma$ 在 $z=0,\infty$ 上方诸极点处的留数形如 $(m_1+m_2,-m_1,-m_2)$，其中 $m_1,m_2\in\C$；而在 $z=q,1$ 上方，它们形如 $(2m,-m,-m)$，其中 $m\in\C$。这表明满穿孔对应于 SU(3) 味因子（其复化 Cartan 子代数中的元素为 $(m_1+m_2,-m_1,-m_2)$），而简单穿孔对应于 U(1) 味因子。按定义 21.24 的意义，满穿孔与简单穿孔都是正则的 SU(3) 穿孔。
\end{remark}

\subsection{广义 SU(3) 箭图}

图 21.18 左边线性箭图的 Seiberg–Witten 曲线可用 21.3 节的方法计算，由此可知该理论
% [原文 p.353]
自然地对应于 UV 曲线的某个特定退化极限，在此情形中该 UV 曲线是带五个穿孔的黎曼球面，其中两个是满穿孔、三个是简单穿孔。如 21.4 节那样，我们先假定两个 SU(3) 规范因子都极弱耦合。在 $\tau_2\to1$ 的极限（其中 $\tau_2$ 是第二个节点的规范耦合）下，第二个 SU(3) 规范因子在 Argyres–Seiberg 对偶框架中得到更好的描述，给出图 21.18 右边的箭图，相应的 $S_{0,5}$ 退化画在图的下方。

\Fig[0.9\linewidth]{figures/fig_21_18.png}{21.18}{更一般的 SU(3) 箭图及其相应的 $S_{0,5}$ 退化。}

在这个阶段，不能像 21.4 节那样继续进行：与 $E_6$ SCFT 相耦合的 SU(3) 规范因子的 S 对偶描述是未知的。不过，与 21.4 节类似，可以研究所有可能的广义 SU(3) 箭图。于是可以论证，这些广义 SU(3) 箭图涵盖了下面将引入的理论 $\mathcal{S}[A_2,S_{g,n},D]$ 的所有可能弱耦合描述。广义 SU(3) 箭图的基本构件如下：
\begin{itemize}
\item 九个超多重态的块：
\begin{itemize}
\item 当它们由某个 SU(3) 规范因子的基本表示中的三个超多重态组成时，味对称性为 $\SU(3)\times\U(1)$。
\item 当它们由乘积 $\SU(3)\times\SU(3)$ 的双基本表示中的单个超多重态组成时，味对称性退化为 U(1)。
\end{itemize}
这样的块表示为与两个 SU(3) 因子和一个 U(1) 因子相连的三叉顶点。
\item $E_6$ SCFT，附带 $E_6$ 味对称群的一个 $\SU(3)^3$ 子群的选取。这样的 $E_6$ SCFT 表示为与三个 SU(3) 因子相连的三叉顶点。
\end{itemize}
给定一对 SU(3) 味因子，可以规范化其乘积的对角 SU(3) 子群，从而得到一个 SU(3) 规范群。此外，还可以规范化一个 SU(3) 味因子的 SU(2) 子群，并使之与一个基本超多重态耦合。

% [原文 p.354]
广义 SU(3) 箭图的更多例子见图 21.19。一般地，若箭图中有 $m_2$ 个超多重态块、$m_3$ 个 $E_6$ SCFT、$n_3$ 个 SU(3) 规范群与 $n_2$ 个 SU(2) 规范群，则有 $n_2+n_3$ 个（精确边缘）规范耦合。箭图的显式味对称性由 $f_3=3m_3+2m_2-2n_3-n_2$ 个 SU(3) 因子与 $f_1=m_2+2n_2$ 个 U(1) 因子构成。该箭图有 $g=n_3+1-m_3-m_2$ 个回路。

\Fig[0.95\linewidth]{figures/fig_21_19.png}{21.19}{广义 SU(3) 箭图的例子。}

满穿孔可以由曲面退化产生，也可以由简单穿孔的碰撞产生。例如考虑图 21.20 第一行中的广义 SU(3) 箭图。与这一线性箭图相关联的 Seiberg–Witten 曲线可用 21.3 节的方法计算。在此情形下可以发现，UV 曲线是带六个简单穿孔的黎曼球面。为方便起见，这里把带 $f_1$ 个简单穿孔与 $f_3$ 个满穿孔的黎曼球面记作 $C_{0,(f_1,f_3)}$。$C_{0,(6,0)}$ 的一个可能的退化极限涉及让 SU(3) 节点解耦，如图 21.20 第二行所示。在此极限下，$C_{0,(6,0)}$ 变成两份 $C_{0,(3,1)}$ 的不交并，因为人们得到两个解耦的箭图，各自带一个额外的 SU(3) 味因子。另一个退化发生在例如左侧 SU(2) 规范群解耦时。这恢复了完整的 SU(3) 味对称性，而且对应于两个简单穿孔的碰撞。由此我们看到，两个简单穿孔相撞会产生一个满穿孔，如图 21.20 第三行所示。然后可以通过让相应的 SU(2) 规范群解耦（其中之一是隐含的，只在 SU(3) 规范群的 Argyres–Seiberg 对偶相位中显式出现）来把剩余的简单穿孔两两配对。由此我们断定，Minahan–Nemeschansky $E_6$ SCFT 的 UV 曲线是带三个满穿孔的黎曼球面。

\subsection{六维定义}

对 [366] 中研究的 SU(3) $\mathcal{N}=2$ 箭图理论作 M 理论提升，会涉及一张缠绕 Seiberg–Witten 曲线 $\Sigma$ 的 M5 膜，其中 $\Sigma$ 是某个 UV 曲线 $C_{\mathrm{UV}}$ 的 3 次分歧覆盖，而 $C_{\mathrm{UV}}$ 或是带穿孔的黎曼
% [原文 p.355]
\Fig[0.75\linewidth]{figures/fig_21_20.png}{21.20}{广义 SU(3) 箭图的不同解耦极限。}
球面（对线性箭图），或是带穿孔的椭圆曲线（对项链箭图）。这是如下构型的可归一化形变：三张 M5 膜缠绕 $C_{\mathrm{UV}}$，并与定域在 $C_{\mathrm{UV}}$ 某点处的横向 M5 膜相交。三张 M5 膜构成的多重膜上的世界体积理论是 $A_2$ 型的 $6d$ $\mathcal{N}=(2,0)$ SCFT。在足够低的能量下，横向 M5 膜表现为 $6d$ 理论中定域在 $C_{\mathrm{UV}}$ 中某些点、并沿 $x^{0,1,2,3}$ 延伸的余维 2 缺陷。$A_2$ 型的 $6d$ $\mathcal{N}=(2,0)$ SCFT 允许简单缺陷与满缺陷，它们分别对应于前面描述的简单穿孔与满穿孔。

这一框架可推广到任意亏格 $g$ 与穿孔数 $n$ 的 UV 曲线 $C_{\mathrm{UV}}=C_{g,n}$。$C_{\mathrm{UV}}$ 上的每个穿孔或是简单的、或是满的。把 $6d$ $\mathcal{N}=(2,0)$ 理论在带指定穿孔缺陷的 $C_{\mathrm{UV}}$ 上紧化，再作一次以保持四维 $\mathcal{N}=2$ 超对称为目的的扭曲，所得四维极限给出一个 $4d$ $\mathcal{N}=2$ 超共形场论。$6d$ $\mathcal{N}=(2,0)$ SCFT 允许 $R$-荷为 2 与 3 的保护算子，它们在扭曲后分别变为 $C_{\mathrm{UV}}$ 上的二次与三次微分。所得理论的精确边缘形变空间等同于泰希米勒空间 $T(S_{g,n})$，其中 $S_{g,n}$ 是 $C_{g,n}$ 底下的光滑曲面。对偶群等同于映射类群 $\mathrm{Mod}(S_{g,n})$。

% [原文 p.356]
\begin{definition}[21.34] 以这种方式得到的 $4d$ $\mathcal{N}=2$ 理论称为 $A_2$ 型 S 类理论，记作 $\mathcal{S}[A_2,C_{g,n},\mathcal{D}]$，其中 $C_{g,n}$ 是 UV 曲线，$\mathcal{D}$ 是缺陷数据，即 $C_{g,n}$ 的哪些穿孔是简单的、哪些是满的。这类理论的库仑分支对应于 $C_{g,n}$ 上二次微分与三次微分 $\phi_2$ 与 $\phi_3$ 构成的空间。在无质量形变时，$\phi_2$ 在每个穿孔处有 1 阶极点，而 $\phi_3$ 在简单穿孔处有 1 阶极点、在满穿孔处有 2 阶极点。Seiberg–Witten（SW）曲线是 $C_{g,n}$ 在 $T^*_{1,0}C_{g,n}$ 中的一个 3 次分歧覆盖，由方程
\[
\lambda^3=\phi_2\lambda+\phi_3
\]
确定，其中 $\lambda$ 表示 $T^*_{1,0}C_{g,n}$ 上的全纯 Liouville 1-形式。精确边缘形变的空间是 $T(S_{g,n})$，对偶群是 $\mathrm{Mod}(S_{g,n})$。
\end{definition}

\begin{remark}[21.35] 也可以考虑带非正则穿孔的 S 类理论。这类理论同样记作 $\mathcal{S}[A_2,C_{g,n},\mathcal{D}]$，其中 $\mathcal{D}$ 编码每个穿孔的精确信息。这允许描述一大类渐近自由理论。与 $A_1$ 型相反，在 $A_2$ 型中还可以考虑另一种穿孔，这与紧李代数 $\mathfrak{su}(3)$ 具有非平凡外自同构群 $\Z_2$ [338] 有关。这类穿孔涉及按该外自同构作扭曲，因此称为扭曲穿孔（twisted punctures）。
\end{remark}

\section{$A_{N-1}$ 型 S 类理论}

$A_{N-1}$ 型 S 类理论的构造与 $A_2$ 型非常相似，区别在于正则穿孔的分类除简单穿孔与满穿孔外还包括新的情形。

带 $2N$ 个无质量基本超多重态的 SU($N$) SYM 的 Seiberg–Witten 曲线可用 21.3 节的方法得到。人们得到方程：
\[
\lambda^N=\sum_{i=2}^{N}\frac{u^{(i)}\dd z^i}{(z-1)(z-q)z^{i-1}}\,\lambda^{N-i}
=\sum_{i=2}^{N}\phi_i(z)\lambda^{N-i},
\]
其中 $\lambda$ 是 $T^*_{1,0}\mathbb{CP}^1_z$ 上的 Liouville 1-形式，每个 $\phi_i$ 是 $\mathbb{CP}^1_z$ 上在穿孔 $z=0,q,1,\infty$ 之外全纯的 $i$-微分。每个 $\phi_i$ 在 $z=q,1$ 处有 1 阶极点，在 $z=0,\infty$ 处有 $i-1$ 阶极点。

\begin{definition}[21.36] 使所有微分 $\phi_i$ 都有 1 阶极点的穿孔称为简单穿孔；使每个 $\phi_i$ 都有 $i-1$ 阶极点的穿孔称为满穿孔。简单穿孔与满穿孔都是正则的。
\end{definition}

如前所述，带 $2N$ 个基本味的 SU($N$) SYM 的 UV 曲线是带两个简单穿孔与两个满穿孔的黎曼球面 $\mathbb{CP}^1$。

% [原文 p.357]
当基本超多重态被赋予一般质量时，可以显式验证：$\lambda$ 在 0 与 $\infty$ 上方诸极点处的留数形如 $(m_1+\dots+m_{N-1},-m_1,\dots,-m_{N-1})$。这个向量属于 $\mathfrak{su}(N)$ 的（复化）Cartan 子代数。与此同时，在 $q$ 与 $1$ 上方诸极点处的留数形如 $((N-1)m,-m,\dots,-m)$。这个向量属于 $\mathfrak{u}(1)$ 的（复化）Cartan 子代数，其中 $\mathfrak{u}(1)\simeq\C$。因此，人们给每个简单穿孔联系一个 U(1) 味因子，给每个满穿孔联系一个 SU($N$) 味因子。

图 21.21 所示的线性共形箭图中，圆节点对应于 SU($N$) 规范因子，方节点对应于 U($N$) 味因子；它由 $k$ 个带双基本超多重态的 SU($N$) 规范群组成，并带第一个与最后一个 SU($N$) 规范因子的 $N$ 个基本表示。该箭图的显式味对称性 $\SU(N)\times\U(1)^{k+1}\times\SU(N)$ 编码在理论的 UV 曲线中，该曲线是带 $k+1$ 个简单穿孔与 2 个满穿孔的黎曼球面。

\Fig[0.95\linewidth]{figures/fig_21_21.png}{21.21}{SU($N$) 规范群的线性共形箭图。}

与前面各节类似，$A_{N-1}$ 型 S 类理论如下定义：

\begin{definition}[21.37] $A_{N-1}$ 型 S 类理论是 $A_{N-1}$ 型 $6d$ $\mathcal{N}=(2,0)$ 理论（即 $N$ 张重合 NS5 膜的世界体积理论）在称为 UV 曲线的带穿孔黎曼曲面 $C_{g,n}$ 上紧化、并作部分扭曲以在四维保持 $4d$ $\mathcal{N}=2$ 超对称所得的四维极限。$A_{N-1}$ 型的 $6d$ $\mathcal{N}=(2,0)$ SCFT 有 $R$-荷为 $2,3,\dots,N$ 的保护算子，它们在扭曲后分别变为 $C_{g,n}$ 上阶为 $2,3,\dots,N$ 的全纯微分 $\phi_2,\phi_3,\dots,\phi_N$。这些全纯微分在穿孔处的行为编码在缺陷数据 $\mathcal{D}$ 中：给定这些行为相当于在 $6d$ 理论中插入余维 2 缺陷，它们定域在 $C_{g,n}$ 的穿孔处，并沿 $x^{0,1,2,3}$ 延伸。该理论的 Seiberg–Witten 曲线是 $C_{g,n}$ 在 $T^*_{1,0}C_{g,n}$ 中的一个 $N$ 次分歧覆盖，由方程
\begin{equation}
\lambda^N=\sum_i\phi_i\lambda^{N-i}, \tag{21.16}
\end{equation}
确定，其中 $\lambda$ 表示 $T^*_{1,0}C_{g,n}$ 上的全纯 Liouville 1-形式。$A_{N-1}$ 型 S 类理论记作 $\mathcal{S}[A_{N-1},C_{g,n},\mathcal{D}]$。对超共形理论，精确边缘形变的空间是 $T(S_{g,n})$，对偶群是 $\mathrm{Mod}(S_{g,n})$，其中 $S_{g,n}$ 是 $C_{g,n}$ 底下的光滑曲面。
\end{definition}

% [原文 p.358]
\begin{remark}[21.38] 在 $A_{N-1}$ 型 S 类中构造 $4d$ $\mathcal{N}=2$ SCFT 的一种标准做法，是只考虑定义 21.24 所定义的正则穿孔。然而，如同 $A_2$ 型，穿孔也可能是非正则的或扭曲的。$A_{N-1}$ 型 S 类理论的非正则穿孔与扭曲穿孔已在 [76, 79] 中研究。
\end{remark}

暂且不论非正则穿孔与扭曲穿孔的情形，我们现在回顾 $A_{N-1}$ 型中正则穿孔的分类。

\subsection{正则穿孔}

\begin{proposition}[21.39] 正则穿孔可以刻画为这样的穿孔：在无质量形变时，对所有的 $i$，全纯微分 $\phi_i$ 在该处的极点阶数至多为 $i-1$。当存在质量形变时，微分 $\phi_i$ 可以有一阶为 $i$ 的极点，从而保证微分 $\lambda$ 在覆盖映射 $\Sigma\to C_{g,n}$ 下穿孔的原像处有简单极点，且留数对应于质量参数，其中 $\Sigma$ 是由式 (21.16) 定义的 Seiberg–Witten 曲线。
\end{proposition}

简单穿孔与满穿孔代表正则穿孔的极端情形。然而当 $N>3$ 时，中间情形也成为可能。对每个 $i$，设 $p_i<i$ 表示在无质量形变时微分 $\phi_i$ 在给定穿孔处的极点阶数。在 $C_{g,n}$ 的某个穿孔附近，覆盖 $\Sigma\to C_{g,n}$ 的各分支分组为若干块，每块包含行为相同的分支。于是每个穿孔由一串整数 $(q_1,\dots,q_s)$ 刻画，其中每个 $q_i$ 表示一块中分支的数目。这串整数满足：
\[
\sum_s q_s=N,
\]
这表明正则穿孔对应于 $N$ 的分拆：不失一般性，可以假定 $(q_1,\dots,q_s)$ 是非增序列。

\begin{definition}[21.40] $N$ 的一个分拆是一个非增整数序列 $q=(q_1,\dots,q_s)$，使得
\[
\sum_{i=1}^{s}q_i=N.
\]
\end{definition}

分拆的另一种常见记号如下。若 $q_1=\dots=q_{i_1}>q_{i_1+1}=\dots=q_{i_1+i_2}>q_{i_1+i_2+1}\dots$，则令 $\chi_1=q_1$、$\chi_2=q_{i_1+1}$ 等等，并写作
\[
q=\chi_1^{i_1}\chi_2^{i_2}\dots.
\]
$N$ 的每个分拆都可以用一个杨图（Young diagram）表示，它由 $N$ 个方格组成，按长度 $q_1,q_2,\dots,q_s$ 从左对齐、自下而上排列。共有 $i_1$ 行长度为 $\chi_1=q_1$，$i_2$ 行长度为 $\chi_2=q_{i_1+1}$，依此类推。

% [原文 p.359]
若杨图包含 $j_1$ 行长度为 $\psi_1$、$j_2$ 行长度为 $\psi_2$ 等等，则与 $q$ 共轭的分拆是 $q^T=\psi_1^{j_1}\psi_2^{j_2}\dots$，它也是 $N$ 的一个分拆。$N=15$ 的三个分拆及其共轭如图 21.22 所示。

\Fig[0.95\linewidth]{figures/fig_21_22.png}{21.22}{$15$ 的三个分拆及其共轭。}

在存在质量形变时，穿孔由 $\lambda$ 在覆盖映射 $\Sigma\to C_{g,n}$ 下该穿孔原像处的留数集合的形式来刻画。正如上文中简单穿孔与满穿孔的例子所示，并非所有留数都必然不同。相反，留数可以按每块 $h_t$ 个、且全部留数之和为零的方式分组，从而产生 $N$ 的另一个分拆 $h$。

$q$、$h$ 与无质量形变时的极点阶数 $p_i$ 之间的联系，在 [146, 第 4 节] 中有详细说明：

\begin{proposition}[21.41] 在无质量形变时，设 $q$ 表示刻画某个穿孔的分拆。则在有质量时，有 $h=q^T$。此外，若把 $h$ 写作 $h=\psi_1^{j_1}\psi_2^{j_2}\dots$，则与该穿孔相关联的味对称性是
\[
G_q=\mathrm{S}\left(\prod_\alpha \U(j_\alpha)\right).
\]
事实上，对应于分拆 $h$ 的留数集合可以自然地看作属于 $G_q$ 的复化 Cartan 子代数。最后，为继续构造，把数字 1 到 $N$ 填入对应于 $q$ 的杨图的方格中，从最左下角的方格开始填 1，沿行从左向右填，然后逐列向上移动。设 $h(i)$ 表示含数字 $i$ 的方格的高度。则对每个 $i$，$p_i=i-h_i$。注意总有 $p_1=0$——注意式 (21.16) 中没有 1 阶全纯微分——并且 $p_i\le i-1$。这表明所考虑的穿孔按命题 21.39 是正则穿孔。
\end{proposition}

\begin{example}[21.42] 对 $N=3$，可能的分拆为：
\begin{itemize}
\item $q=1^3$，$h=3^1$，对应于三个相等、和为零的留数，即只有零留数，且 $p_1=p_2=p_3=0$。这一情形对应于没有穿孔，而不是一个穿孔。
% [原文 p.360]
\item $q=2^11^1$，$h=2^11^1$，对应于留数集合 $\{2m,-m,-m\}$，且 $p_2=p_3=1$。这是简单穿孔的情形，味对称性为 $\mathrm{S}(\U(1)\times\U(1))=\U(1)$。
\item $q=3^1$，$h=1^3$，对应于留数集合 $\{m_1+m_2,-m_1,-m_2\}$，$p_2=1$，$p_3=2$。这是满穿孔的情形，相应的味对称性为 SU(3)。
\end{itemize}
\end{example}

\begin{example}[21.43] 对 $N=4$，可能的分拆为：
\begin{itemize}
\item $q=1^4$，对应于无奇点。
\item $q=2^11^2$，$h=3^11^1$，对应于留数 $\{3m,-m,-m,-m\}$，且 $p_2=p_3=p_4=1$。这是简单穿孔的情形，味对称性为 $\mathrm{S}(\U(1)\times\U(1))\simeq\U(1)$。
\item $q=2^2$，$h=2^2$，对应于留数 $\{m,m,-m,-m\}$，$p_2=p_3=1$，且 $p_3=2$。味对称性为 SU(2)。
\item $q=3^11^1$，$h=2^11^2$，对应于留数 $\{2m_1+m_2,-m_1,-m_1,-m_2\}$，$p_2=1$，且 $p_3=p_4=2$。味对称性为 $\mathrm{S}(\U(1)\times\U(2))\simeq\U(1)\times\SU(2)$。
\item $q=4^1$，$h=1^4$，对应于留数 $\{m_1+m_2+m_3,-m_1,-m_2,-m_3\}$，$p_2=1$，$p_3=2$，$p_4=3$。这是满穿孔的情形，味对称性为 SU(4)。
\end{itemize}
\end{example}

因此，在只含正则穿孔的理论 $\mathcal{S}[A_{N-1},C_{g,n},\mathcal{D}]$ 中，缺陷数据 $\mathcal{D}$ 只在于给 UV 曲线 $C_{g,n}$ 上的每个穿孔指定一个含 $N$ 个方格的杨图。在无质量形变时，这个杨图编码全纯微分 $\phi_2,\dots,\phi_N$ 在每个穿孔处的极点阶数。一般质量形变的形式也由 $\mathcal{D}$ 规定，如上所述。注意当 $N=3$ 时，正则穿孔必为简单穿孔或满穿孔，这与定义 21.34 一致。

\subsection{退化极限}

当 UV 曲线 $C_{g,n}$ 的亏格 $g$ 为 0 或 1 时，带一般正则穿孔的 $C_{g,n}$ 的某些退化极限对应于理论 $\mathcal{S}[A_{N-1},C_{g,n},\mathcal{D}]$ 的弱耦合线性或项链共形箭图描述。这类箭图的 Seiberg–Witten 曲线可用 21.3 节的方法计算。

例如，考虑规范群为 $\SU(N_1)\times\SU(N_2)\times\dots\times\SU(N_k)$ 的一般线性共形箭图，如图 21.4 所示，其中 $\max(N_i)=N$。令 $N_0=N_{k+1}=0$，并设 $F_i$ 表示与每个规范因子 $\SU(N_i)$ 相关联的基本超多重态数目，$i=1,\dots,k$。该箭图理论是共形的，只要对所有的 $1\le i\le k$ 都有
\[
2N_i-N_{i-1}-N_{i+1}=F_i\ge0.
\]
这一要求意味着序列 $(N_1,\dots,N_k)$ 是凹的，从而满足：

% [原文 p.361]
\[
0<N_1<\dots<N_{k_1}=N=N_{k_1+1}=\dots=N_{k_2}>\dots>N_k>0,
\]
其中 $1\le k_1\le k_2\le k$ 为某些下标。对每个 $1\le i\le k-1$，定义 $q^L_i:=N_i-N_{i-1}$，并令
\[
q^L=(q^L_1,\dots,q^L_{k_1}).
\]
$(N_1,\dots,N_k)$ 的凹性意味着 $q^L$ 是 $N$ 的一个分拆。类似地，存在对应于箭图右端的分拆 $q^R$。按 21.3 节的思路计算这类箭图的 Seiberg–Witten 曲线表明，UV 曲线是带穿孔的黎曼球面 $C_{0,n}$，其中 $n-2$ 个穿孔是简单穿孔，其余两个由箭图左右两端所决定的分拆 $q^R$ 与 $q^L$ 刻画。$N=4$ 的两个例子如图 21.23 所示。

\Fig[0.8\linewidth]{figures/fig_21_23.png}{21.23}{$A_3$ 型 S 类理论的弱耦合极限。}

一般地，$\mathcal{S}[A_{N-1},C_{g,n},\mathcal{D}]$ 的弱耦合描述可以表述为广义 SU($N$) 箭图。这些广义 SU($N$) 箭图的构件中包括本质强耦合的 $4d$ $\mathcal{N}=2$ SCFT，它们推广了广义 SU(3) 箭图中出现的 $E_6$ 理论的情形。$A_{N-1}$ 型中可用的正则穿孔种类繁多，允许大量这类推广。

例如，带 $2N$ 个基本超多重态的 SU($N$) SYM 的 UV 曲线是带两个简单穿孔与两个满穿孔的黎曼球面。对 $N=3$，我们已经看到两个简单穿孔合并的退化极限对应于带六个基本超多重态的 SU(3) SYM 的 Argyres–Seiberg 对偶相位。更一般地，对任意 $N$，都可以考虑两个简单穿孔相撞的退化极限，由此得到 Argyres–Seiberg 对偶的一种推广：当 $\tau\to1$ 时，带 $2N$ 个基本超多重态的 SU($N$) SYM 更好地描述为一个 SU(2) 规范理论，它与一个基本超多重态以及一个秩为 $N-2$、味对称性为 $\SU(2)\times\SU(2N)$ 的相互作用 $4d$ $\mathcal{N}=2$ SCFT 相耦合。这个相互作用理论有时记作
% [原文 p.362]
$R_N$，其 UV 曲线是带两个满穿孔与一个简单穿孔的球面（参见 [76] 与 [339, 第 12 节]）。

$E_6$ SCFT 的另一种可能推广是 $A_{N-1}$ 型的 S 类理论，其 UV 曲线是带三个满穿孔的黎曼球面。该理论是味对称性为 $\SU(N)^3$ 的相互作用 $4d$ $\mathcal{N}=2$ SCFT，有时称为 $T_N$（参见 [339, 第 12 节]）。

一般正则穿孔相撞的退化极限在 [146] 中研究。一般地，这类碰撞导致弱耦合的 $\SU(M)$ 规范群，其中 $M\le N$。然而有一个特殊情形：当两个穿孔相撞、且它们所对应的杨图都只有两列时，这不会产生弱耦合的特殊酉规范群，而产生弱耦合的 USp 规范群。因此，以 USp 规范群结尾的特殊酉群箭图属于广义 SU($N$) 箭图的范畴。一个具体实例给出 Argyres–Seiberg 对偶的另一个例子：在强耦合极限下，带六个味的 USp(4) 对应于与 Minahan–Nemeschansky $E_7$ SCFT 相耦合的弱耦合 SU(2) 规范理论 [29]。

\section{一般 S 类理论}

具有 $\mathcal{N}=(2,0)$ 超对称的六维超共形场论（SCFT）的存在性，最初是由超弦理论与 M 理论的考虑所论证的 [365, 315, 321, 320]。这些理论具有时空对称代数 $\mathfrak{osp}(6,2|4)$，它是最大可能的超共形代数。已知的 $6d$ $\mathcal{N}=(2,0)$ SCFT 分为两大类：阿贝尔 $6d$ $\mathcal{N}=(2,0)$ SCFT 与相互作用 $6d$ $\mathcal{N}=(2,0)$ SCFT，它们各自与 Dynkin 型为 $A$、$D$ 或 $E$ 的简单紧李代数相关联。与李代数 $\mathfrak{u}(1)$ 相联系的阿贝尔理论描述单张 M5 膜上的世界体积动力学。相反，对应于 $\mathfrak{su}(N)$（$A_{N-1}$ 型）的相互作用理论描述 $N$ 张 M5 膜构成的多重膜的世界体积理论。关于 $6d$ $\mathcal{N}=(2,0)$ SCFT 更广泛的讨论可参见 [194] 及其中的参考文献。

IIA 型超弦理论中 $N$ 张 NS5 膜构成的多重膜 T 对偶于 IIB 型超弦理论中的 $\C^2/\Z_N$ 奇点。六维中 $\mathcal{N}=(2,0)$ 超对称的保持与如下事实有关：$\C^2/\Z_N$ 是仿射 Calabi–Yau 二重奇点，它保持 II 型超弦理论中 32 个超荷的一半。更一般地，保持 32 个超荷中 16 个的奇点，其和乐必须落在 SU(2) 中，这表明它们具有 $\C^2/\Gamma$ 的形式，其中 $\Gamma$ 是 SU(2) 的离散子群。SU(2) 离散子群的分类遵循 $ADE$ 型模式，这反映了这些奇点的类型。通过 T 对偶，这些奇点实现对应于每个 $ADE$ 型的 $6d$ $\mathcal{N}=(2,0)$ SCFT。

$A_{N-1}$ 理论联系着 M5 膜多重膜上的世界体积动力学，它在我们讨论中一直处于核心地位，因为它在带缺陷数据 $\mathcal{D}$ 的带穿孔黎曼曲面 $C_{g,n}$ 上紧化并作部分扭曲后，在四维
% [原文 p.363]
极限下给出 $A_{N-1}$ 型 S 类理论。在这些理论中，Seiberg–Witten 曲线 $\Sigma$ 是 $C_{g,n}$ 在余切丛 $T^*_{(1,0)}C_{g,n}$ 中的全纯分歧覆盖，其中限制后的全纯 Liouville 1-形式在 $\Sigma$ 上定义 Seiberg–Witten 微分。$A_{N-1}$ 型 S 类理论涵盖了许多已知的四维 $\mathcal{N}=2$ 理论，不过也存在一些 $4d$ $\mathcal{N}=2$ 理论，其在 S 类中的实现仍然难以捉摸，甚至可能无法实现。

如先前所提到的，S 类构造的一个重要特征，是它把 Seiberg–Witten 理论与 Hitchin 系统显式地联系起来 [114]。具体地说，四维 $\mathcal{N}=2$ 理论的库仑分支是一个全纯可积系统的底，而 Hitchin 系统提供了关键的例子。尽管整个可积系统结构——即库仑分支上的环面纤维化——对 $4d$ 理论而言可能显得是辅助性的，但当理论进一步在圆上紧化时，它就获得了物理意义。在此情形下，该可积系统直接对应于所得三维理论的库仑分支（例如参见 [322, 151]）。

特别地，对 $A_{N-1}$ 型 S 类理论，编码其低能库仑分支物理的可积系统是对应于 $G_{\C}=\SL_n(\C)$ 的 Hitchin 系统。这一结构只有在圆上紧化之后才完全显现。从 $A_{N-1}$ 型的六维 $\mathcal{N}=(2,0)$ 理论出发，这相当于先在带缺陷数据 $\mathcal{D}$ 的 $C_{g,n}$ 上紧化，再在 $S^1$ 上紧化。或者，若先把 $6d$ 理论在 $S^1$ 上紧化，则得到规范群为 $A_{N-1}$ 的五维 $\mathcal{N}=2$ SYM 理论。再在 $C_{g,n}$ 上紧化时，这个 $5d$ SYM 理论的 BPS 方程恰好约化为 Hitchin 方程。这就解释了为什么 Hitchin 系统作为对应于 $A_{N-1}$ 型 S 类理论的可积系统出现 [155]。更一般地，这表明 S 类理论恰好就是其 Seiberg–Witten 可积系统为 Hitchin 系统的那些 $4d$ $\mathcal{N}=2$ 理论。

同样的构造也适用于从 $D$ 型或 $E$ 型的一般相互作用 $6d$ $\mathcal{N}=(2,0)$ 理论出发，得到 $D$ 型或 $E$ 型的 S 类理论，同样记作 $\mathcal{S}[\g,C_{g,n},\mathcal{D}]$（其中 $\g$ 为 $D$ 型或 $E$ 型），它们对应于 $D$ 型或 $E$ 型的 Hitchin 系统。这些理论允许以规范因子为 $D$ 型或 $E$ 型的规范理论的形式给出弱耦合描述。这些类型的 S 类理论中的穿孔也可以是正则的、非正则的或扭曲的 [77, 78, 81, 80, 82, 83, 84, 85]。

与非单边群（即 $B$、$C$、$F$、$G$ 系列）相关联的 Hitchin 系统，来自把与 $A$、$D$ 或 $E$ 型李代数 $\g$ 相关的 $6d$ $\mathcal{N}=(2,0)$ SCFT 在圆上、带 $\g$ 的一个外自同构扭曲地紧化，从而得到 $B$、$C$、$F$、$G$ 型的 $\mathcal{N}=2$ $5d$ SYM 理论 [202], [203], [230]。

关于 S 类理论近期进展的综述，读者可参阅综述文章 [9] 及其所引文献。

% [原文 p.364]
% 原书此页为空白页（第 21 章正文结束于印刷 p.363 / PDF p.403）。

% [Tutor Note: 图 21.11 的记号约定：圆节点（内部）是 SU(2) 规范因子，方节点是 SU(2) 味因子；双基本超多重态必然出现在一条内部边上，其两个端点各带一个规范群，因此图 21.11 左图在一个规范节点上多加一个味节点后，每个 SU(2) 规范因子才具有四个基本味（两个显式 + 两个来自双基本）。核对弱耦合描述个数：带 5 个标记点的三叉树（即 S_{0,5} 的裤子分解）共 15 棵，恰与命题 21.25 中 15 个尖点相符；计数公式为 5!/(2^2·2!)=15。]

% [Tutor Note: p.345 的外边数公式可做量纲/欧拉公式核对：连通三叉树有 V 个顶点与 E_i 条内部边，则 3V=2E_i+E_e，故 E_e=3V-2E_i；又树满足 V=E_i+1，代入得 E_e=3(E_i+1)-2E_i=E_i+3=n+3。取 n=1（两个三叉顶点由一条内部边相连），四个外部边确实围成 S_{0,4} 的一个分解。]

% [Tutor Note: p.349 的同余子群中原文写作 $c\equiv0\ [2]$（方括号表示模 2），与常见的 $c\equiv0\ (\mathrm{mod}\ 2)$ 同义；$\Gamma_0(2)$ 是 $\PSL_2(\Z)$ 的指数 3 子群，故 $\HH/\Gamma_0(2)$ 有两个尖点，这正是文中 $\tau\to1$ 与 $\tau\to\infty$ 不等价的几何来源。]

% [Tutor Note: p.351 的曲线 (Σ) 中，$\phi_2$ 在 $0,q,1,\infty$ 各有一阶极点、$\phi_3$ 在 $q,1$ 一阶而在 $0,\infty$ 二阶；这与定义 21.32 中「$\phi_2$ 一阶、$\phi_3$ 二阶 ⟺ 满穿孔」一致：满穿孔 $b,d$ 在 $0,\infty$，简单穿孔 $a,c$ 在 $q,1$。核对时可把 (Σ) 写成 $\lambda^3-\phi_2(z)\lambda-\phi_3(z)=0$，再按 $\lambda=x\,\dd z/z$ 比较系数：$\phi_2\propto\dd z^2/((z-1)(z-q)z)$ 在有限处为 0,q,1 三个一阶极点、在 $z=\infty$ 亦为一阶极点；$\phi_3\propto\dd z^3/((z-1)(z-q)z^2)$ 在 0,∞ 为二阶、在 q,1 为一阶。]

% [Tutor Note: p.359 的公式 $p_i=i-h(i)$ 可用来核对两个例子：对 $q=(2,2)$（Example 21.43 第三项），杨图两行各 2 格，按「自左下起、沿行向右、再逐列向上」填号得 $h(1)=h(2)=1$、$h(3)=h(4)=2$，故 $p_2=1,p_3=1,p_4=2$。原文此处末句印作「$p_3=2$」，应为 $p_4=2$（原文笔误，照录并在此指出）；其余各项与本公式一致。]

% [Tutor Note: p.354 的构件计数可做一致性核对：取纯 $E_6$ SCFT 构件的箭图（$m_3=1$，$m_2=n_3=n_2=0$），公式给出 $f_3=3$ 个 SU(3) 因子、$f_1=0$、回路数 $g=1+0-1-0=0$，正与「一个 $E_6$ SCFT 作为连接三个 SU(3) 因子的三叉顶点」相符（对照图 21.20 最后一行的 $C_{0,(0,3)}$）。再取纯九个超多重态块（$m_2=1$，其余为 0），得 $f_3=2$、$f_1=1$、$g=0$，也与「三叉顶点连接两个 SU(3) 与一个 U(1)」相符。]

% [Tutor Note: (21.16) 中原文把求和写成 $\sum_i$ 而未标出上下限；对照 p.356 的同一曲线（显式写成 $\sum_{i=2}^{N}$）以及定义 21.37 中 $\phi_2,\dots,\phi_N$ 的约定，此处的 $i$ 应在 $2\le i\le N$ 范围内求和。译文照录原文的 $\sum_i$，读者按 $i=2,\dots,N$ 理解即可。]
